\section{Relation between Euclidean Steiner Tree and MMX}

\subsection{Euclidean Steiner Tree}

Given a set of points $P$ in $\mathbb{R}^N$, the Euclidean Steiner Tree problem is to find the shortest tree (possibly
adding new points) that connects all points in $P$. Calling the resulting graph $G = (V, E)$, the cost of the tree is
\begin{equation}
    \sum_{(u, v) \in E} \|u - v\|
\end{equation}
Where $\| \cdot \|$ is the Euclidean norm.

\subsection{MMX}
The MMX is defined in \ref{Paper linkato} as follows: given a set of points $P$ in $\mathcal{R}^N$, the minimization problem is:

\begin{align}
    \text{(MMX)}: \; \min &\sum_{[i,j] \in E_1} \|a^i - x^j\| y_{ij} + \sum_{[i,j] \in E_2} \|x^i - x^j\| y_{ij}, \\
    & \sum_{j \in S} y_{ij} = 1 \quad \text{for } i \in P, \\
    & \sum_{i \in P} y_{ij} + \sum_{k < j, k \in S} y_{kj} + \sum_{k > j, k \in S} y_{jk} = 3 \quad \text{for } j \in S, \label{MMX:connectivity_steiner} \\
    & \sum_{k < j, k \in S} y_{kj} = 1 \quad \text{for } j \in S - \{p+1\},  \\
    & y_{ij} \in \{0, 1\}, \quad [i, j] \in E, \quad x^i \in \mathbb{R}^n, \quad i \in S.
\end{align}

Where: \begin{itemize}
           \item $P$ is the set of terminals,
           \item $S$ is the set of Steiner points (in this case $S = \{|P| + 1, ..., 2|P|-2\}$),
           \item $E_1$ is the set of edges connecting terminals to Steiner points,
           \item $E_2$ is the set of edges connecting Steiner points
\end{itemize}

\subsection{Relation between the two problems}

The MMX problem "solves" the ESTP in the sense that under certain conditions, the optimal solution of MMX corresponds to
the optimal EST.
Explicitly, we can construct a function $\phi$ that maps feasible solutions of MMX to feasible solutions of ESTP and such
that the following commutative diagram holds:

\[
    \begin{tikzcd}
        & \mathbb{R}  \\
        MMX \arrow[ur, "Obj"] \arrow[rr, "\phi"'] & & EST  \arrow[ul, "Obj"]
    \end{tikzcd}
\]

Where $Obj$ is the objective function of the respective problems and MMX and EST are the feasible regions.

\subsection{Definition of $\phi$}
The function $\phi$ maps $\mathbb{R}^{N (|P|-2)} \times \{0,1\}^{|E|}$ to a graph in $\mathbb{R}^N$ as follows:

\begin{enumerate}
    \item For each $i \in P$, add a terminal at $a^i$.
    \item For each $j \in S$, add a vertex $v_j$ at $x^j$.
    \item For each $[i,j] \in E$ with $y_{ij} = 1$, add an edge between $v_i$ and $v_j$.
\end{enumerate}

There are a few issues with this identification: $\phi$ is neither injective (there is a permutation invariance in the index of both
Steiner points and terminals) nor surjective, as for instance, the number of vertices is fixed. It is known \ref{Boh} that the number of Steiner points
in an optimal EST cannot be bigger than the ones included in MMX, but it could be lower (in the paper it is specified that
this formulation only holds for full Steiner topologies).
It is easy to see that the objective functions are the same as they both sum the Euclidean distances of the active edges.


\textbf{Claim} Calling MMX and EST the respective feasible regions, we have that $\phi(MMX) \subseteq EST$.

\begin{proof}
    Let $x \in MMX$ and $G = \phi(x)$. We need to show that $G$ is a tree (so connected and loop-free).
    By the first constraints, we know each terminal is connected to a Steiner point and constraint (\ref{MMX:connectivity_steiner}) ensures that
    all Steiner points are connected, so $G$ is connected.
    By edge counting, we deduce that $G$ is a tree.
\end{proof}



\textbf{Claim} The optimal solution of MMX is also optimal for ESTP.

\begin{proof}
    We know from \ref{Libro di optimal transport} that the optimal solution for the ESTP is a tree with
    number of vertices $|V| \le 2|P| - 2$ and degree of each vertex $d(v) \le 3$.
    We can show that for each graph with those properties there is an element in $MMX$ that maps to it.
    In $MMX$ the degree of terminals is fixed to 1 but we can add a Steiner point with coordinates equal to the terminal
    (so allowing a degenerate solution) and make the following construction (for the case $d(v) = 2$):

    OLD PART:
    We add a vertex $i$ with coordinate $x_i = a_j$, we activate edge (of length 0) connecting $i$ to $j$. Then the two
    edges connecting $j$ to the other points are attached to $i$, in this way the degree of the degenerate Steiner point is
    3 and the degree of the terminal is 1.
    We can do a similar construction (but with 2 Steiner points) for the case $d(v) = 3$.
    We still need to do some edges and vertices counting to finish the proof (maybe an induction on the number of vertices?).

    NEW PART:
    We will not do an induction but provide a procedure to add Steiner points until there are enough of them, while keeping
    the graph structure the same. We will then show that those operation preserve all the properties of the MMX problem.

    By doing some edge counting we know that
    \begin{equation}
        |E| = \frac{1}{2} \left( 3|S| + \sum_{i \in P} d(i) \right)
    \end{equation}
    Which implies that:
    \begin{equation}
        \sum_{i \in P} d(i) = 2|P| - 2 - |S|
    \end{equation}
    If $|S| = 2|P| - 2$ we are done, otherwise we can find a terminal $i$ with $d(i) = 2,3$ and add 1 or 2 Steiner points respectively.
    \begin{itemize}

        \item Case $d(i) = 1 $: we add a Steiner point $s$ with connection to the two points connected to $i$ and remove the edge
        connecting $i$ to the other terminal. We also connect $s$ to the terminal $i$.

        \item Case $d(i) = 2$: Calling the points connected to $i$ $z_1, z_2, z_3$ we add two Steiner points $s_1, s_2$ and connect
        $s_1$ to $(z_1, s_2, i)$ and $s_2$ to $(s_1, z_2, z_3)$.
    \end{itemize}

    In this way we preserve connectedness, loop-freeness and the number of (non-degenerate) vertices. Connectedness implies
    \ref{MMX:connectivity_steiner} up to reordering of the Steiner points.
    We can repeat the procedure until $|S| = |P| - 2$ which is our desired number of Steiner points.
    We get that the degree of every terminal is 1 and the degree of every Steiner point is 3, so all the conditions
    of MMX are satisfied.


\end{proof}


\section{Relation between Euclidean Steiner Tree and branched transport}

\subsection{Definition of branched transport}
For now, I will use the Lagrangian formulation of branched transport, as defined in \ref{Libro di branched transport}.

Given $X \subset \mathbb{R}^N$ compact and two measures of equal mass $\mu^+$ and $\mu^-$ on $X$, a traffic plan is
a measure $P$ on $K$ that transports $\mu^+$ to $\mu^-$.
Formally, consider the maps $\pi_0, \pi_\infty : K \rightarrow X$ given by $\pi_0(\gamma) = \gamma(0)$ and $\pi_\infty(\gamma) = \gamma(T(\gamma))$,
$P$ is a traffic plan if:
\[
    \mu^+ = (\pi_0)_{\#}P, \quad \mu^- = (\pi_\infty)_{\#}P
\]

\subsection{Parameterized Traffic Plans}
The next paragraph is copied from the book:


According to Skorokhod theorem (Theorem A.3 p. 185), we can parameterize any traffic plan $P$ by a measurable function
$\chi : \Omega = [0, |\Omega|] \rightarrow K$ such that $P = \chi_{\#} \lambda$, where $\lambda$ is the Lebesgue measure
on $[0, |\Omega|]$. We shall set $\chi(\omega, t) := \chi(\omega)(t)$ and consider it as a function of the variable pair
$(\omega, t)$.

With this formulation the energy (or cost) of a parametrized traffic plan is:
\begin{equation}
    E^{\alpha}(P) = \int_{\mathbb{R}^N} |x|_P^{\alpha} d \mathcal{H}^1(x)
\end{equation}
There are also other formulations and under some assumptions they are equivalent. In particular there should be no problem
for finite atomic measures.

\subsection{Relation between branched transport and ESTP}

We can reduce any ESTP (in the algorithmic sense) to an instance of a branched transport problem.
Given a set of terminals $P = \{a^1, ..., a^{|P|}\}$ we can define:
\begin{equation}
    \mu^+ = \delta_{a^1} \quad \mu^- = \frac{1}{|P| - 1}\sum_{i = 2}^{|P|} \delta_{a^i}
\end{equation}
and set $\alpha =0$ in the cost function.
We can then use results from branched transport to say that the optimal parametrized
traffic plan is a graph and it is loop-free. We will now prove that it is connected and conclude that it is a tree.

For each terminal $a^i$ (with $i>1$) we have a subset $A_i \subset [0,1]$ such that $\chi(A_i)$ are all paths terminating at $a^i$ and
$\chi(A_i)(0) = a_1$.
We can easily conclude that every terminal is connected to $a_1$ and therefore the graph is connected.

